\section{Tools for the finite Gram search}\label{sec:tools}
All bounds below are necessary conditions. We may enlarge a candidate family when doing so makes enumeration simpler; an exclusion of the enlarged family remains valid. We never identify row and column Gram supports without justification.

\subsection{Reduction of the color partitions}\label{sec:partitions}
\begin{lemma}\label{lem:partitions}
In a strict improvement, the row and column color partitions each belong to
\[
 (15),\quad(14,1),\quad(13,2),\quad(13,1,1).
\]
\end{lemma}
\begin{proof}
Suppose a same-color class of size $a$ is internally orthogonal. Its Gram block is $15I_a$. The Schur complement on the other $15-a$ coordinates has trace at most $(15-a)(15-a/5)$, because the cross energy is at least $3a(15-a)$. Thus
\[
 \det G\le15^a(15-a/5)^{15-a}<B\qquad(3\le a\le12).
\]
The displayed inequalities are rational. Consequently, each color class of size between 3 and 12 has internal energy at least 9 in a strict improvement. Listing the partitions subject to this fact and $Q\le168$ leaves the four stated partitions and
\[
 (12,3),\ (12,2,1),\ (11,4),\ (11,3,1),\ (11,2,2),\ (10,5).
\]
For the latter six, use a largest-class block $A$ of size $a$, internal energy $e$, and eigenvalue upper bound $L$. If its cross energy is at least $c$, Schur complementation and the arithmetic--geometric mean inequality give
\begin{equation}\label{eq:partition-bound}
 \det G\le\U_a(e)\left(15-\frac{c}{L(15-a)}\right)^{15-a}.
\end{equation}
The complete list of parameters needed is in \cref{tab:partition}. Every resulting bound is less than $B$. The small-energy eigenvalue caps follow by listing the supports with at most four edges; for the remaining cases the trace--variance eigenvalue bound suffices. The entry catalogue and this finite list are regenerated by \file{proof/order15_mu3_unified_audit/programs/verify_foundations.py}.
\end{proof}
\begin{table}[ht]
\centering
\caption{Parameters in \eqref{eq:partition-bound}; entries in each row are paired in order.}\label{tab:partition}
\begin{tabular}{@{}cccc@{}}\toprule
$a$ & $e$ & $L$ & $c$\\\midrule
10 & $9$ & $18$ &150\\
11 & $9,18,27$ & $20,21,23$ &132\\
12 & $9,18,27,36,45,54$ & $18,77/4,21,22,241/10,25$ &108\\\bottomrule
\end{tabular}
\end{table}

\subsection{Componentwise Schur bounds}
Consider a majority-color block $A=\bigoplus_i A_i$, where component $i$ has $v_i$ vertices and internal energy $9u_i$. Let $D_i\ge\det A_i$ and $L_i\ge\lambda_{\max}(A_i)$. Isolated vertices are included as components with $(v_i,u_i,D_i,L_i)=(1,0,15,15)$.

\begin{lemma}\label{lem:schur}
For a partition $(14,1)$ with total cross energy $c$, set
\[
 s_{14}=\sum_i\frac{3v_i}{L_i}+\frac{c-42}{\max_iL_i}.
\]
Then
\begin{equation}\label{eq:schur14}
 \det G\le\left(\prod_iD_i\right)\max\{0,15-s_{14}\}.
\end{equation}
For a size-13 majority block and two remaining coordinates, let $c$ be the energy from the majority to those coordinates and let $h$ be the norm of their mutual inner product. Put
\[
 s_{13}=\sum_i\frac{6v_i}{L_i}+\frac{c-78}{\max_iL_i}.
\]
If $s_{13}<30$, then
\begin{equation}\label{eq:schur13}
 \det G\le\left(\prod_iD_i\right)
 \left[(15-s_{13}/2)^2-
       \max\{0,\sqrt h-s_{13}/2\}^2\right].
\end{equation}
If $s_{13}\ge30$, no positive definite candidate is possible.
\end{lemma}
\begin{proof}
For one remaining coordinate, write $G=\left(\begin{smallmatrix}A&z\\z^*&15\end{smallmatrix}\right)$. Each coordinate of $z$ has squared modulus at least 3. Therefore
$z^*A^{-1}z\ge\sum_i\|z_i\|^2/L_i\ge s_{14}$, proving \eqref{eq:schur14}.

For two coordinates, write $U^*A^{-1}U=\left(\begin{smallmatrix}x&z\\\bar z&y\end{smallmatrix}\right)$ and $s=x+y$. Then $s\ge s_{13}$ and $|z|\le\sqrt{xy}\le s/2$. The determinant of the remaining Schur block is at most
\[
 (15-s/2)^2-\max\{0,\sqrt h-s/2\}^2.
\]
This expression is decreasing for $0\le s\le30$, since $\sqrt h\le15$ for a positive semidefinite Gram. Its substitution at $s_{13}$ gives the result. A positive definite Schur block must have positive trace $30-s$.
\end{proof}
These estimates retain the loss coming from every isolate. Discarding isolates, or assuming that one can choose the orientation so that they disappear, would invalidate the coverage of several boundary cases.

\subsection{Component determinant bounds}\label{sec:moments}
For each component we take the least of the applicable bounds: the sharp envelope, an exact small-support determinant maximum, and a fourth-moment refinement. For a tree with unit edge norms after division by 3, phases can be switched away for the purpose of computing its spectrum. All nonisomorphic trees of the required order are listed. For connected supports with at most seven vertices and energy at most ten units, all graph-atlas supports, norm assignments, and chord phases are evaluated exactly.

For larger components the following bound avoids an unnecessarily large phase enumeration.
\begin{lemma}\label{lem:moment}
Let $A=15I_v+E\succ0$, with $E$ Hermitian, zero diagonal, $\tr(E^2)=18u$, $\rank E\le d$, and $\lambda_{\max}(A)\le L$. If
\[
 W\ge\sum_{\{i,j,k\}}\frac{|E_{ij}E_{jk}E_{ki}|}{27},
\]
where the sum is over support triangles, then
\begin{equation}\label{eq:moment}
 \log\frac{\det A}{15^v}
 \le-\frac u{25}+\frac{2W}{125}-\frac{3u^2}{125Ld}.
\end{equation}
\end{lemma}
\begin{proof}
For $-1<x\le\rho$, $\rho\ge0$, integration of
$1/(1+x)=1-x+x^2-x^3/(1+x)$ gives
\[
 \log(1+x)\le x-x^2/2+x^3/3-\frac{x^4}{4(1+\rho)}.
\]
Apply this to the eigenvalues of $E/15$, with $\rho=(L-15)/15$. Their first moment is zero. Also $\tr E^3\le162W$ and $\tr E^4\ge(18u)^2/d$. Substitution gives \eqref{eq:moment}.
\end{proof}
A valid choice of eigenvalue cap is
\[
 L=15+\sqrt{18u(v-1)/v}.
\]
If the support has clique number at most $k$, the sharper cap $15+\sqrt{18u(k-1)/k}$ is valid. Indeed, Cauchy--Schwarz bounds $|x^*Ex|^2$ by $36u\sum_{ij\in E(\Gamma)}|x_i|^2|x_j|^2$. For nonnegative weights summing to one, the latter edge sum is at most $(k-1)/(2k)$: repeatedly move the mass between a nonadjacent pair onto one vertex without decreasing the sum, until the support of the weights is a clique, then use Cauchy--Schwarz again.

The triangle bound is also finite and explicit. For a connected graph with $m$ edges and $v$ vertices, put $g=m-v+1$. The number of triangles is bounded using its cycle rank and biconnected blocks. Deleting a minimum-degree vertex of degree $d$ from a 2-connected block reduces the cycle rank to $g-d+1$ and removes at most $\binom d2$ triangles. Together with additivity over blocks, this gives a recursion for an upper bound $T(g)$. The code also takes the minimum with $\binom v3$, $m(m-1)/6$, $g(g+1)/2$, and $2^{g-1}$. The last bound counts the vectors of odd edge parity in the cycle space: there are either none or exactly $2^{g-1}$, and every triangle gives one. For a norm assignment $w_e=|E_e|^2/9$, put $\Delta_e=\sqrt{w_e}-1$. Expanding each triangle weight gives the valid upper bound
\[
 W\le T(g)+g\sum_e\Delta_e+
       \sum_{e<f}\Delta_e\Delta_f+
       \sum_{e<f<h}\Delta_e\Delta_f\Delta_h.
\]
An edge belongs to at most $g$ triangles, and two edges to at most one, which justifies the coefficients. The alternative $W\le T(g)\sqrt{w_{(1)}w_{(2)}w_{(3)}}$, using the three largest weights, is also applied. For each weight partition the smaller of these two bounds is used; maximizing this value over all allowed weight partitions supplies a bound depending only on $(v,u)$.

All square roots in these comparisons are enclosed by rational bounds with the direction of rounding chosen explicitly. When \eqref{eq:moment} gives $\det A\le15^v e^{-a}$ with $a>0$, the program uses
\[
 e^{-a}\le\left(\sum_{j=0}^{24}\frac{a^j}{j!}\right)^{-1}.
\]
No decimal approximation is used to discard a candidate. If a rank or eigenvalue cap is known, the moment optimization is repeated with that cap: enumerate the number of eigenvalues at the cap and the two possible remaining stationary values. Boundary points with determinant zero need not be retained.

\subsection{Two-sided rank and spectral restrictions}
For pure-color Grams put $T=(G-15I)/3$ and $S=(K-15I)/3$. They satisfy
\begin{equation}\label{eq:intertwine}
 TH=HS.
\end{equation}
\begin{lemma}[Isolates and kernel projections]\label{lem:projection}
If the support of $T$ has $r$ isolates, then
\[
 15(P_{\ker S})_{jj}\ge r\quad\text{for every }j.
\]
The analogous statement holds with $T$ and $S$ exchanged. In particular, an isolate on one side forbids a support leaf on the other side.
\end{lemma}
\begin{proof}
The $r$ corresponding rows of $H$ are mutually orthogonal, each has squared norm 15, and by \eqref{eq:intertwine} lies in $\ker S$. If they form $R$, then $R^*R/15\preceq P_{\ker S}$. Every diagonal entry of $R^*R$ equals $r$. At a support leaf, the kernel equation forces the coordinate at its neighbor to vanish, contradicting a positive diagonal lower bound there.
\end{proof}
The inequalities must be applied in both directions, matching characteristic polynomials rather than assuming equal support graphs. If $r_{\min}$ is a lower bound for the isolate count on the opposite side, summing the projection inequality over a component of size $v$ gives nullity at least $\lceil vr_{\min}/15\rceil$. Iterating this restriction with the determinant bounds often removes all component configurations.

\begin{lemma}[Mixed-color rank signatures]\label{lem:mixed-rank}
Suppose both color partitions are $(14,1)$ and both majority supports have isolates. On the row side let $r$ be the isolate count, $k$ the nullity of the nonisolated majority block after subtracting $15I$, and $t$ the number of coordinates on which its kernel can be nonzero. Use $R,K,T$ for the corresponding column-side quantities. Then
\begin{equation}\label{eq:mixed-signature}
 r+k=R+K,\qquad r(15-R-T)\le15,\qquad
 r(15-R)\le15(K+1),
\end{equation}
and the two inequalities with the sides exchanged also hold.
\end{lemma}
\begin{proof}
In the kernel of $G-15I$, an isolated-majority row forces the minority coordinate to be zero. The nonisolated coordinates then belong to a $k$-dimensional kernel, and the minority equation imposes one nontrivial relation on those coordinates and the $r$ isolate coordinates. Thus the nullity is $r+k-1$. Matching nullities gives the first equation.

Let $Y$ consist of the $r$ isolated-majority rows of $H$. Then $YY^*=15I_r$ and $Y(K-15I)$ has rank at most one. Modulo one common row direction, its rows lie in $\ker(K-15I)$. That kernel is zero outside the $R$ isolate coordinates and the $T$ kernel-support coordinates. Restricting $Y$ to the other $15-R-T$ columns therefore gives rank at most one. Its squared Frobenius norm is $r(15-R-T)$, while its squared operator norm is at most 15. This proves the second inequality. After deleting only the $R$ isolate columns, the rank is at most $K+1$, proving the third in the same way.
\end{proof}

\begin{theorem}[Spectral rectangle inequality]\label{thm:rectangle}
Suppose $HH^*=nI+aT$ and $H^*H=nI+aS$, where $T,S$ are Hermitian, $a>0$, and every entry of $H$ has modulus one. Let $I,J$ be unions of invariant coordinate blocks of $T,S$. If
\[
 p(x)=\gcd(\chi_{T_I}(x),\chi_{S_J}(x))
      =x^d+c_1x^{d-1}+\cdots,
\]
then
\begin{equation}\label{eq:rectangle}
 |I||J|\le nd-ac_1.
\end{equation}
For $d=0$, the right side is interpreted as zero.
\end{theorem}
\begin{proof}
Put $X=H_{I,J}$. Invariance gives $T_IX=XS_J$, so the spectral decomposition of $X$ has nonzero blocks only at common eigenvalues. At a common eigenvalue $\lambda$, the rank is at most the smaller multiplicity and every squared singular value is at most $n+a\lambda$, because $XX^*\preceq nI+aT_I$. Summing these bounds with minimum multiplicities gives $nd-ac_1$. On the other hand, $\|X\|_F^2=|I||J|$.
\end{proof}
This test uses the trace of the common spectral part, not merely its dimension or the largest eigenvalue of the whole block. In our application $n=15$ and $a=3$.
